\subsection{README}

This is the TU Delft starterkit for open publishing with Jupyter Book. It provides a skeleton with basic settings to have a head start with your thesis / project.

It includes:

\begin{itemize}
\item a github deploy file taking care of deploying the website and building a pdf
\item basic content files that can be altered to fit your project
\item yml files in the root folder to configure the book and the pdf build
\end{itemize}

\subsubsection{🕵 Where to start}

Already know about Jupyter Book, familiar with GitHub or GitLab, Markdown... start with cloning the \href{https://github.com/new?template\_name=starterkit\&template\_owner=TUD-JB-OS}{starterkit} and start writing working on your project. New to the ecosystem? Start with \href{https://tud-jb-os.github.io/book}{the manual}, and then move on to the starterkit.

\subsubsection{🏁 Purpose}

The purpose of this project and this manual is to enable others to use Jupyter Book for writing and publishing their own scientific and educational content. We hope to lower the technical barrier for researchers and educators to create and share their own resources, and to promote open science and open education practices.

\subsubsection{📄 License}

This book is licensed under the \href{http://creativecommons.org/licenses/by-nc/4.0/}{Creative Commons Attribution-NonCommercial 4.0 International License}, unless stated otherwise.

You are free to:

\begin{itemize}
\item Share --- copy and redistribute the material in any medium or format
\item Adapt --- remix, transform, and build upon the material for any non-commercial purpose provided proper attribution is given.
\end{itemize}

\subsubsection{❌ Errors}

If you find any errors in the content, please report them by opening an issue on the \href{https://github.com/TUD-JB-OS/book/issues}{GitHub repository}.

\subsubsection{💰 Funding}

This project has received funding from the \href{https://www.tudelft.nl/en/open-science/articles-tu-delft/mainstreaming-open-science-fund-2025-2026}{Delft University of Technology Open Science Fund (2025-2026)}.

\subsubsection{⭐ Identifier}

TUDJBOSSTARTERKIT